% ========= PRAMBEL =========
% Hier werden die Grundeinstellungen des Dokuments festgelegt und Pakete geladen.

\documentclass[
    12pt,           % Schriftgröße
    a4paper,        % Papierformat
    oneside         % Einseitiger Druck (alternativ: twoside für beidseitigen Druck)
]{scrartcl}         % KOMA-Script Klasse, gut für europäische Standards

% --- KODIERUNG UND SPRACHE ---
% Diese Pakete sind für deutschen Text unerlässlich.
\usepackage[utf8]{inputenc}   % Erlaubt die direkte Eingabe von Umlauten (ä, ö, ü) und ß im Editor.
\usepackage[T1]{fontenc}      % Stellt sicher, dass Schriften korrekt kodiert werden, was für die Silbentrennung von Wörtern mit Umlauten wichtig ist.
\usepackage[ngerman]{babel}   % Lädt deutsche Sprachanpassungen (z.B. "Inhaltsverzeichnis" statt "Contents", korrekte Trennregeln).

% --- NÜTZLICHE PAKETE ---
\usepackage{amsmath}          % Umfassende Erweiterungen für mathematische Formeln.
\usepackage{amssymb}          % Stellt zusätzliche mathematische Symbole bereit.
\usepackage{graphicx}         % Ermöglicht das Einbinden von Grafiken (JPG, PNG, PDF).
\usepackage{geometry}         % Zum einfachen Einstellen der Seitenränder.
\usepackage{hyperref}         % Erstellt klickbare Links im PDF (z.B. im Inhaltsverzeichnis oder bei Querverweisen).

% --- EINSTELLUNGEN FÜR PAKETE ---
\geometry{
    a4paper,
    top=2.5cm,
    bottom=2.5cm,
    left=3cm,
    right=3cm
}

\hypersetup{
    colorlinks=true,          % Links farbig statt mit Rahmen
    linkcolor=blue,           % Farbe für interne Links (Inhaltsverzeichnis, etc.)
    citecolor=green,          % Farbe für Zitat-Links (mit biblatex/bibtex)
    urlcolor=magenta          % Farbe für externe URLs
}

% --- DOKUMENTINFORMATIONEN ---
% Diese Informationen werden für die Titelseite verwendet.
\title{Titel der Arbeit}
\author{Max Mustermann}
\date{\today} % Setzt das heutige Datum automatisch. Alternativ: \date{15. Oktober 2023}

% ========= DOKUMENT =========
% Der eigentliche Inhalt des Dokuments beginnt hier.
\begin{document}

% --- TITELSEITE ---
% Erstellt die Titelseite basierend auf den obigen \title, \author und \date Befehlen.
\maketitle

% --- INHALTSVERZEICHNIS ---
% Erstellt automatisch ein Inhaltsverzeichnis aus den \section, \subsection, etc. Befehlen.
\tableofcontents
\newpage % Beginnt nach dem Inhaltsverzeichnis eine neue Seite.

% --- KAPITEL/ABSCHNITTE ---
\section{Einleitung}
Dies ist ein Beispieltext für eine LaTeX-Vorlage auf Deutsch. Hier können Umlaute wie ä, ö, ü und auch das scharfe S (ß) problemlos verwendet werden.
Dank des \texttt{babel}-Pakets mit der Option \texttt{ngerman} funktioniert die deutsche Silbentrennung automatisch und korrekt, auch bei schwierigen Wörtern wie Donaudampfschifffahrtsgesellschaftskapitän.

Man kann auch Fußnoten\footnote{Dies ist eine Fußnote.} ganz einfach einfügen.

\subsection{Motivation}
Hier steht die Motivation für dieses Dokument. Man kann auch auf andere Abschnitte verweisen, zum Beispiel auf Abschnitt~\ref{sec:beispiele}. Das Tilde-Zeichen (~) sorgt für ein geschütztes Leerzeichen, damit die Nummer nicht in die nächste Zeile rutscht.

\section{Weitere Beispiele}
\label{sec:beispiele} % Eine Marke, auf die man sich beziehen kann.

Hier folgen einige typische Elemente, die in wissenschaftlichen Arbeiten oder Dokumenten vorkommen.

\subsection{Listen}
Man kann sowohl ungeordnete als auch geordnete Listen erstellen.

\begin{itemize}
    \item Erster Punkt einer ungeordneten Liste.
    \item Zweiter Punkt mit einem längeren Text, der bei Bedarf automatisch umgebrochen wird.
    \item Und so weiter...
\end{itemize}

\begin{enumerate}
    \item Erster Punkt einer geordneten Liste.
    \item Zweiter Punkt.
    \item Dritter Punkt.
\end{enumerate}

\subsection{Mathematische Formeln}
LaTeX ist besonders stark im Satz von mathematischen Formeln. Man kann sie direkt im Text einfügen, wie z.B. die berühmte Formel $E = mc^2$, oder als abgesetzte, nummerierte Formel:
\begin{equation}
    \int_{a}^{b} x^2 \,dx = \frac{b^3 - a^3}{3}
    \label{eq:integral}
\end{equation}
Auf diese Formel (\ref{eq:integral}) kann man sich im Text beziehen.

\subsection{Tabellen}
Eine einfache Tabelle könnte so aussehen:

\begin{center} % Zentriert die Tabelle auf der Seite
\begin{tabular}{|l|c|r|}
    \hline
    \textbf{Spalte 1 (links)} & \textbf{Spalte 2 (zentriert)} & \textbf{Spalte 3 (rechts)} \
    \hline
    Zeile 1, Zelle 1 & Zelle 2 & 123.45 \
    Zeile 2, Zelle 1 & Zelle 2 & 6.7 \
    \hline
\end{tabular}
\end{center}

\section{Zusammenfassung}
Hier endet das Beispieldokument. Viel Erfolg bei der Arbeit mit LaTeX!

\end{document}
